\subsection{The exact first integral kernel of the unramified comparison}

\begin{statement}{E.102}{Proved}{All finite factors and inverse maps in the degree-four full-to-metabelian fibre}
Retain the actual cubic geometric pro-three group $N$,
its intrinsic metabelian quotient $M=N/N''$, and the
original arithmetic automorphisms $T,J$ of E.100.
For the induced degree-four lattice
\[
 K_4=(N''\cap\gamma_4N)/(N''\cap\gamma_5N)
       =\ker(L_4(N)\longrightarrow L_4(M)),
 \tag{E102.1}
\]
the complete integral coefficient quotient is
\[
 \mathcal C_4(N'')=K_4/(T_4-I)K_4^J
   \simeq\mathbb Z_3^{504}\oplus C_3^{144}
       \oplus C_9^{56}\oplus C_{27}^{30}\oplus C_{81}^{9}.
 \tag{E102.2}
\]
Here $C_m=\mathbb Z/m\mathbb Z$ with its finite topology.
The finite subgroup has order $3^{382}$ and exponent $81$.
All isomorphisms below are continuous homomorphisms of
compact $\mathbb Z_3$-modules, and their complete integral
basis maps and inverses are retained.

This computes the entire first nonzero coefficient kernel
in E100.12. Through the actual triangular comparison
E100.24--E100.26, it gives every degree-four fibre above
the full unramified metabelian section space, including
the zero-inertia loci in E.98 and E.100.
The higher-degree factors keep their exact presentations
and inverse maps of E.100. The manuscript anchor is
\S51.16, folios733v--736v, scans172--178,
formulas(371)--(379), Part E(c),(d) and follow-up4.
\end{statement}

\begin{proof}
\emph{The actual integer map.}
We use the ordered non-Chen Hall basis of E.100.
It identifies $K_4$ with $\mathbb Z_3^{990}$; the
integral arithmetic matrices $T_4,J_4$ in this very basis
are the saved restrictions of the actual Hall actions.
In particular this basis retains the original marking,
without replacing the arithmetic operators by chosen
operators with the same ranks. The complete matrices
satisfy
\[
 J_4^2=I,\qquad \operatorname{tr}(J_4)=-18,
 \qquad\operatorname{rank}K_4^J=486.
 \tag{E102.3}
\]
The identity $J_4^2=I$ gives the direct decomposition
by the two integral projectors $(I\pm J_4)/2$.
Therefore $\operatorname{im}(I+J_4)=K_4^J$ and
\[
 C=(T_4-I)(I+J_4):\mathbb Z_3^{990}\longrightarrow K_4
 \quad\hbox{has image }(T_4-I)K_4^J.
 \tag{E102.4}
\]
The fixed rank in (E102.3) gives
$\operatorname{rank}_{\mathbb Q_3}C\le486$.

\emph{Finite exact computation and its full integral lift.}
The retained program performs integer row and column
operations modulo $q=3^7=2187$. At each step it takes
an entry of minimal three-adic valuation, moves it to
the next pivot, multiplies its row by the inverse of
its unit part, and clears that column and row.
All coefficients and both inverse transformations
are recorded. The result is
\[
 UCV\equiv\operatorname{diag}(D,0_{504})\pmod q,
 \qquad
 D=\operatorname{diag}
       (1^{247},3^{144},9^{56},27^{30},81^9).
 \tag{E102.5}
\]
The certificate verifies both products of $U$ and its
retained modular inverse, both products of $V$ and its
retained modular inverse, and every entry of (E102.5).
These give $486$ nonzero pivots modulo $q$, so the
upper rank bound is attained:
$\operatorname{rank}_{\mathbb Q_3}C=486$.
In particular no remaining nonzero invariant factor
is hidden by the finite precision.

Here is an explicit lift of these maps over
$\mathbb Z_3$. Choose the integer representatives
between zero and $q-1$ of the saved matrices $U,V$.
They are invertible over $\mathbb Z_3$: their displayed
modular inverse identities already give invertibility
modulo three. Write their exact integer product as
\[
 B=UCV=\begin{pmatrix}B_{11}&B_{12}\\B_{21}&B_{22}\end{pmatrix},
 \qquad B_{11}=D S,\quad
 S=I+E,\quad E=D^{-1}(B_{11}-D).
 \tag{E102.6}
\]
Every entry of $B-\operatorname{diag}(D,0)$ is divisible
by $3^7$. Every diagonal entry of $D$ divides $3^4$.
Thus $E$ is integral and divisible by $3^3$, and
\[
 S^{-1}=\sum_{j\ge0}(-E)^j
 \tag{E102.7}
\]
converges in the integral matrix ring.
Both $S^{-1}D^{-1}B_{12}$ and
$B_{21}S^{-1}D^{-1}$ are integral.
Consequently the explicit triangular matrices
\[
 \begin{aligned}
 V_1&=\begin{pmatrix}
 S^{-1}&-S^{-1}D^{-1}B_{12}\\0&I
 \end{pmatrix},\\
 U_1&=\begin{pmatrix}
 I&0\\-B_{21}S^{-1}D^{-1}&I
 \end{pmatrix}
 \end{aligned}
 \tag{E102.8}
\]
and their block-triangular inverses belong to
$\mathrm{GL}_{990}(\mathbb Z_3)$.
Their exact product is
\[
 U_1 B V_1=
 \begin{pmatrix}
 D&0\\0&B_{22}-B_{21}S^{-1}D^{-1}B_{12}
 \end{pmatrix}
       =\operatorname{diag}(D,0_{504}).
 \tag{E102.9}
\]
The final lower block vanishes over $\mathbb Q_3$
because the full matrix has rank $486$ and its
upper block already has that rank. All its entries
are integral, so it vanishes over $\mathbb Z_3$ as well.
This proves the complete integral Smith equality,
with the actual matrices $U_1U$ and $VV_1$.

For completeness, inverses of the integer lifts $U,V$
can also be recovered directly from the retained
modular inverse matrices. Let $\widetilde U^{-1}$
denote the stored integer representative and put
$E_U=I-U\widetilde U^{-1}$. Then
\[
 U^{-1}=\widetilde U^{-1}
             \sum_{j\ge0}E_U^j,
       \qquad E_U\in3^7M_{990}(\mathbb Z_3),
 \tag{E102.10}
\]
and the same formula gives $V^{-1}$.
Every series in these formulas has a stated integral
convergence rate and determines all digits of its map.

\emph{Quotient coordinates and the section-space comparison.}
For $x\in K_4$, apply $U_1U$ to its coordinate column.
The first $247$ coordinates have quotient zero;
reduce the next $144,56,30,9$ coordinates modulo
$3,9,27,81$, respectively, and retain the last $504$
coordinates in $\mathbb Z_3$.
Equation (E102.9) proves that the exact kernel of this
continuous homomorphism is $(T_4-I)K_4^J$.
To recover a preimage, choose the standard integer
representatives in the finite coordinates, insert
zeros in the first $247$ positions, and apply
$U^{-1}U_1^{-1}$. This is a continuous set section;
on the quotient it gives the inverse module
isomorphism, independent of the chosen finite lifts.
These are the two maps in (E102.2).
Their torsion order is
$3^{144+2\cdot56+3\cdot30+4\cdot9}=3^{382}$,
and the nine $C_{81}$ factors prove the exponent claim.

The strict kernel and fixed-lift maps in E100.11--12
identify this quotient with the entire kernel of
$\mathcal C_4(N)\to\mathcal C_4(M)$.
All earlier induced kernel layers vanish.
The nonlinear aligning conjugator and arithmetic
correction at the next step are exactly E100.24--25;
inserting the representatives just constructed into
E100.26 gives every degree-four preimage.
Thus the calculation transfers to the stated section
fibres by those actual maps. The entire inverse limit
continues with E.100's subsequent integral coordinates.

The files \texttt{e\_kernel\_degree4\_smith.py},
its actual stdout, its JSON operation log, and its
NumPy archive retain every matrix in (E102.3)--(E102.6),
all modular inverses and the original Hall indices.
All arithmetic is integer arithmetic. The largest
absolute entry of $C$ is $4800$, and the exact
integer-lift product $B$ is checked with explicit
overflow bounds before each machine-integer product.
The complete identities, rather than only the Smith
list, are the finite certificate for (E102.5).
\end{proof}
